% ============================================================
% TC3002B (ACA-2026) — Proyecto Triton Heal (Equipo 11)
% Poster: evaluación experimental de generación de kernels
% Triton con SLM + verificación pre-GPU y método dual.
% Contenido tomado de report/main.pdf (rúbrica 1.1–2.7).
% Convención: español para comunicar, inglés para código.
% ============================================================

\documentclass[final]{beamer}

% ====================
% Packages
% ====================
\usepackage[T1]{fontenc}
\usepackage[utf8]{luainputenc}
\usepackage{lmodern}
\usepackage[size=custom, width=122,height=91, scale=1.2]{beamerposter}
\usetheme{gemini}
\usecolortheme{msu}
\usepackage{graphicx}
\usepackage{booktabs}
\usepackage{tikz}
\usetikzlibrary{arrows.meta, positioning, shapes.geometric}
\usepackage{pgfplots}
\pgfplotsset{compat=1.14}
\usepackage{anyfontsize}

% ====================
% Lengths
% ====================
\newlength{\sepwidth}
\newlength{\colwidth}
\setlength{\sepwidth}{0.025\paperwidth}
\setlength{\colwidth}{0.3\paperwidth}

\newcommand{\separatorcolumn}{\begin{column}{\sepwidth}\end{column}}

% ====================
% Title
% ====================
\title{Triton Heal: generación de kernels Triton con SLM locales y verificación pre-GPU como filtro de seguridad}

\author{Cristóbal Camarena Hernández \and Josué Galindo Gutiérrez \and Gandhi Emmanuel Valdez Huerta \and Diego Alberto Estrada López}

\institute[shortinst]{Tecnológico de Monterrey --- TC3002B (ACA-2026) --- Equipo 11}

% ====================
% Footer
% ====================
\footercontent{
  \href{https://github.com/G4ndhiV/TRITONHEAL}{Repo: github.com/G4ndhiV/TRITONHEAL} \hfill
  ACA-2026 --- Sesión con socio formador \hfill
  \href{mailto:a01642653@tec.mx}{a01642653@tec.mx}}

% ====================
% Logo (cabecera)
% ====================
\logoleft{\tikz\node[fill=white, rounded corners=10pt, inner sep=12pt]
  {
\includegraphics[height=6cm]{logos/logo.png}};}

% ====================
% Body
% ====================
\begin{document}

\begin{frame}[t]
\begin{columns}[t]
\separatorcolumn

% ============================================================
% COLUMNA 1 — Motivación, problema, antecedentes, diferenciador
% ============================================================
\begin{column}{\colwidth}

  \begin{block}{Resumen}

    Los modelos de lenguaje pequeños (SLM) producen kernels en Triton DSL con
    errores que solo se descubren al compilar en GPU, un recurso caro.
    \textbf{Triton Heal} propone un pipeline que \emph{genera y luego verifica
    antes de tocar la GPU}: un SLM local (Llama-3.1-8B) escribe el kernel y, si
    no pasa las validaciones, un modelo de código (DeepSeek-R1-Distill) lo repara.
    Evaluamos seis configuraciones sobre \textbf{10 tareas de generación}
    (50 inferencias) y un banco de \textbf{70 kernels} de verificación
    (420 evaluaciones) en un \textbf{Mac M4} con Ollama, con diseño
    \emph{within-subjects} y pruebas pareadas (McNemar / Wilcoxon / Friedman,
    Holm). DeepSeek-Coder-V2 logró el mejor balance (\textbf{100\% compilación,
    MSR 87.1\%, FNR 12.9\%}); el dual alcanzó \textbf{MSR 80.6\% / F1 0.86},
    mejorando la detección de Llama solo sin superar al mejor frontera.
    Reportamos con honestidad estadística: la mayoría de contrastes \emph{no}
    rechaza $H_0$ tras Holm.

  \end{block}

  \begin{block}{1. Planteamiento del problema}

    \begin{itemize}
      \item \textbf{Dolor:} compilar y ejecutar kernels Triton en GPU es lento
        y costoso; un SLM puede generar muchos kernels inválidos que
        desperdician ese tiempo de cómputo.
      \item \textbf{Línea base:} usar un único modelo grande de frontera por
        API es caro y depende de conectividad; un SLM local solo es más barato
        pero falla más y no avisa cuándo se equivoca.
      \item \textbf{Pregunta de investigación:} ¿un esquema \emph{dual}
        (SLM local + reparación con modelo de código) más un filtro de
        verificación previo a la GPU mejora la tasa de compilación frente a
        GPT-4o, y reduce los fallos no detectados (FNR) frente a usar un solo
        SLM?
    \end{itemize}

  \end{block}

  \begin{block}{2. Antecedentes}

    \begin{itemize}
      \item \textbf{Triton DSL:} lenguaje de GPU embebido en Python; permite
        escribir kernels con \texttt{@triton.jit} y operaciones
        \texttt{tl.load}/\texttt{tl.store}, más tratable que CUDA crudo.
      \item \textbf{SLM vs.\ frontera:} los modelos pequeños (8B–14B) corren en
        local (Ollama, Apple Silicon) pero tienen menor tasa de acierto que
        modelos de frontera; el reto es cerrar esa brecha sin pagar API.
      \item \textbf{Verificación pre-compilación:} en vez de compilar a ciegas,
        un \emph{checker} valida sintaxis (AST de Python), reglas léxicas y
        estructura Triton, y clasifica el kernel \texttt{safe}/\texttt{unsafe}
        (métricas MSR, FNR, F1) antes de gastar la GPU.
    \end{itemize}

  \end{block}

  \begin{alertblock}{3. Nuestro diferenciador}

    \textbf{En una línea:} mientras otros enfoques apuestan a un solo modelo
    grande, nosotros combinamos un \emph{SLM local que genera} con un
    \emph{modelo de código que repara} y un \emph{filtro de verificación
    pre-GPU}, evaluado con un diseño experimental pre-registrado.

    \begin{itemize}
      \item \textbf{Método dual (propio).} Llama-3.1-8B genera; si el checker
        rechaza, DeepSeek-R1-Distill lo repara (hasta 2 reintentos) y un veto
        final decide \texttt{safe}/\texttt{unsafe}.
      \item \textbf{Rigor experimental.} 6 configuraciones, $T=0$, semilla 42,
        diseño \emph{within-subjects}, hipótesis pre-registradas (1.1–1.10) y
        corrección Holm-Bonferroni sobre la familia de contrastes.
      \item \textbf{Honestidad de resultados.} Declaramos límites: M4 sin CUDA,
        compilación \emph{sintáctica} (no ejecución GPU), y baselines
        Claude/GPT.
    \end{itemize}

  \end{alertblock}

\end{column}

\separatorcolumn

% ============================================================
% COLUMNA 2 — Arquitectura, diseño experimental, métricas
% ============================================================
\begin{column}{\colwidth}

  \begin{block}{4. Arquitectura del pipeline}

    \begin{figure}
      \centering
      \begin{tikzpicture}[
          node distance=1.1cm and 0.9cm,
          pbox/.style={draw, rounded corners, minimum height=2.0cm,
                       text width=6.0cm, align=center, font=\small, thick},
          pbio/.style={pbox, fill=blue!8},
          pbproc/.style={pbox, fill=green!12},
          pbrep/.style={pbox, fill=red!8},
          pbout/.style={pbox, fill=orange!14},
          arr/.style={-{Stealth[length=4mm]}, very thick}]
        \node[pbio]   (task)  {\textbf{Tarea}\\\texttt{tasks.jsonl}};
        \node[pbproc] (gen)   [right=of task] {\textbf{Generación SLM}\\Llama-3.1-8B (Ollama)};
        \node[pbproc] (chk)   [right=of gen]  {\textbf{Checker L1--L3}\\AST + léxico + \texttt{@triton.jit}};
        \node[pbrep]  (rep)   [below=of chk]  {\textbf{Reparación}\\DeepSeek-R1-Distill ($\leq 2$)};
        \node[pbproc] (ver)   [left=of rep]   {\textbf{Verificación pre-GPU}\\MSR / FNR / F1};
        \node[pbout]  (out)   [left=of ver]   {\textbf{Veredicto}\\\texttt{safe} / \texttt{unsafe}};
        \draw[arr] (task) -- (gen);
        \draw[arr] (gen)  -- (chk);
        \draw[arr] (chk)  -- node[right,font=\footnotesize]{falla} (rep);
        \draw[arr] (rep)  -- (ver);
        \draw[arr] (ver)  -- (out);
        \draw[arr] (chk.south west) .. controls +(-1,-1) .. node[above,font=\footnotesize]{compila} (ver.north);
      \end{tikzpicture}
      \caption{De una tarea a un veredicto: el SLM genera, el checker valida
      en tres niveles, R1-Distill repara los fallos y la verificación pre-GPU
      etiqueta el kernel sin compilar en CUDA.}
    \end{figure}

    El checker exige tres niveles en M4: \textbf{L1} sintaxis Python (AST),
    \textbf{L2} reglas léxicas del equipo, \textbf{L3} presencia de
    \texttt{@triton.jit} con \texttt{tl.load}/\texttt{tl.store}. Compilar aquí
    significa pasar L1--L3, \emph{no} ejecutar en GPU NVIDIA.

  \end{block}

  \begin{block}{5. Diseño experimental}

    Seis configuraciones comparadas con el mismo benchmark y tareas
    (\emph{within-subjects}, $T=0$, semilla 42, \texttt{timeout} 120\,s):

    \begin{itemize}
      \item \textbf{SLM locales:} Llama-3.1-8B, DeepSeek-R1-Distill-Qwen-14B.
      \item \textbf{Frontera:} DeepSeek-Coder-V2, Claude 3.5 Sonnet, GPT-4o
      
      \item \textbf{Propuesto:} Triton Heal dual (Llama + reparación DeepSeek-R1-Distill-Qwen).
    \end{itemize}

    \vspace{0.4em}
    \noindent\textbf{Contrastes pre-registrados (Holm-Bonferroni):}
    {\small
    \begin{table}
      \centering
      \begin{tabular}{l l l}
        \toprule
        \textbf{ID} & \textbf{Prueba} & \textbf{Comparación} \\
        \midrule
        C4 & McNemar  & Compilación: dual vs GPT-4o \\
        C7 & McNemar  & Compilación: dual vs Llama \\
        C1--C3 & McNemar & Detección (MSR/FNR) \\
        C5 & Wilcoxon & Latencia frontera vs SLM \\
        C6 & Friedman & F1 entre configuraciones \\
        \bottomrule
      \end{tabular}
    \end{table}}

  \end{block}

  \begin{block}{6. Métricas de evaluación}

    \begin{itemize}
      \item \textbf{Tasa de compilación:} \% de salidas que pasan L1--L3.
      \item \textbf{Speedup proxy:} $s = t_{\text{baseline}}/t_{\text{pipeline}}$,
        indicador operativo en M4 (no es speedup medido en GPU).
      \item \textbf{MSR} (Mutant-Spotting Rate): fracción de kernels
        \texttt{unsafe} correctamente marcados; \textbf{FNR}: \texttt{unsafe}
        pasados por alto.
      \item \textbf{F1} y \textbf{\% JSON válido}: calidad y parseabilidad de la
        respuesta del verificador.
    \end{itemize}

    Relevancia práctica exigida: $\geq 5$ pp en compilación/MSR y $\sim 8$ pp de
    reducción de FNR, además de $p_{\text{Holm}} < 0.05$.

  \end{block}

\end{column}

\separatorcolumn

% ============================================================
% COLUMNA 3 — Resultados, discusión, referencias
% ============================================================
\begin{column}{\colwidth}

  \begin{block}{7. Resultados}

    \textbf{Titular:} sobre 10 tareas de generación y un banco de 70 kernels
    (Mac M4, Ollama, semilla 42), DeepSeek-Coder-V2 fue el más confiable
    (\textbf{100\% compilación, MSR 87.1\%, FNR 12.9\%}); el dual logró
    \textbf{MSR 80.6\% / F1 0.86}, por encima de Llama solo (64.5\%) en
    detección, sin superar al mejor frontera. Tras Holm, solo el Friedman
    global de F1 es significativo ($p < 0.001$).

    \begin{table}
      \centering
      {\footnotesize\setlength{\tabcolsep}{4pt}
      \begin{tabular}{l r r r r}
        \toprule
        \textbf{Configuración} & \textbf{Comp.\%} & \textbf{MSR\%} & \textbf{F1} & \textbf{FNR\%} \\
        \midrule
        Llama-3.1-8B (local)      & 50  & 64.5 & 0.81 & 35.5 \\
        R1-Distill-14B (local)    & 0   & 12.9 & 0.22 & 87.1 \\
        \textbf{Coder-V2} (front.) & \textbf{100} & \textbf{87.1} & 0.81 & \textbf{12.9} \\
        Claude 3.5 (heur.)        & 100 & 83.9 & 0.91 & 16.1 \\
        GPT-4o (heur.)            & 100 & 83.9 & 0.91 & 16.1 \\
        \textbf{Triton Heal (dual)} & 50 & 80.6 & 0.86 & 19.4 \\
        \bottomrule
      \end{tabular}}
      \caption{Descriptivos por configuración (gen.\ $N=10$, verif.\ $N=70$).}
    \end{table}

    \begin{figure}
      \centering
      \begin{minipage}{0.49\linewidth}\centering
        \includegraphics[width=\linewidth]{figures/fig_compile_rate.pdf}
      \end{minipage}\hfill
      \begin{minipage}{0.49\linewidth}\centering
        \includegraphics[width=\linewidth]{figures/fig_msr_f1.pdf}
      \end{minipage}
      \caption{Izq.: tasa de compilación (\%) de kernels generados.
      Der.: MSR (\%) y F1$\times$100 en verificación.}
    \end{figure}

    \begin{figure}
      \centering
      \includegraphics[width=0.82\textwidth]{figures/fig_robustness_compile.pdf}
      \caption{Mapa de compilación tarea $\times$ configuración
      (verde = éxito, rojo = fallo).}
    \end{figure}

  \end{block}

  \begin{block}{8. Conclusiones y trabajo futuro}

    \begin{itemize}
      \item \textbf{Conclusión:} el mejor desempeño global fue de
        DeepSeek-Coder-V2; el dual es prometedor como \emph{arquitectura}
        (reparación local + veto de frontera), no como ganador absoluto.
      \item \textbf{Evidencia:} tras Holm ningún McNemar rechaza $H_0$ (el dual
        no supera a GPT-4o, $p\approx 0.15$); solo el Friedman global de F1 es
        significativo ($p<0.001$).
      \item \textbf{Limitaciones:} M4 sin CUDA (compilación \emph{sintáctica});
        muestra reducida; R1-14B inestable; baselines Claude/GPT.
      \item \textbf{Trabajo futuro:} APIs reales, 20 tareas $\times$ 3 réplicas,
        GPU NVIDIA y exportar el filtro a decodificación restringida.
    \end{itemize}

  \end{block}

  \begin{block}{Referencias}

    {\scriptsize\raggedright
    [1] P.~Tillet et al. \emph{Triton: An IL and Compiler for Tiled Neural
    Network Computations.} MAPL, 2019.\\[0.2em]
    [2] Q.~McNemar. \emph{Note on the sampling error of correlated proportions.}
    Psychometrika, 1947.\\[0.2em]
    [3] S.~Holm. \emph{A simple sequentially rejective multiple test procedure.}
    Scand.\ J.\ Stat., 1979.\\[0.2em]
    [4] M.~Friedman. \emph{The use of ranks to avoid the normality assumption in
    ANOVA.} JASA, 1937.\par}

  \end{block}

\end{column}

\separatorcolumn
\end{columns}
\end{frame}

\end{document}
